% Hand-written narrative; numbers are macros.
Across all \ExpNAlerts{} alerts the frozen detector emits at 10 alerts per day on the test windows
(\ExpNTrue{} true), the family an explanation names is itself informative (Table~\ref{tab:family}).
Alerts explained by a \emph{volume} deviation are true \ExpVolPrec\% of the time, against
\ExpAllPrec\% over all alerts, and they account for \ExpVolShareTrue\% of the true alerts; no alert
explained by an \emph{identity} deviation was true (\ExpIdnPrec\%), in line with the ablation. An
analyst who works the queue by explanation type rather than by rank would reach true alerts sooner.
% Hand-written narrative; numbers are macros.
As an independent check of faithfulness, Kernel SHAP~\cite{lundberg2017shap} was run on the
detector's score for \ShapN{} alerts of one test window, with each user's usual feature vector as
the baseline. SHAP's top-ranked feature is one of the features the explanation names in
\ShapTopNamed\% of alerts, and it lies in the explanation's primary family in \ShapFamilyAgree\%.


Whether a person can triage from the explanation alone (H3) is tested by a pre-registered crossover
study: \StudyN{} alerts (\StudyTrue{} true, from \StudyInsiders{} insiders across all three
scenarios, and \StudyFalse{} false), two arms
(score only vs.\ explanation only), four reviewers, each alert judged twice per arm by different
people, and a paired bootstrap over alerts on balanced accuracy. The study is \TriageStatus{}
(Table~\ref{tab:triage}); all \TriageNUsable{} of \TriageNExpected{} planned judgements were usable.
Reviewers shown only the explanation reached a balanced accuracy of \TriageBAB{}, against
\TriageBAA{} for reviewers shown only the score, a difference of \TriageDiff{} percentage points
(95\% CI \TriageLo{} to \TriageHi{}). The entire interval lies above the pre-registered
+\TriageThr{} points threshold, so H3 is \textsc{\TriageVerdict}. Reviewers in arm~B were
substantially more accurate on both true alerts (sensitivity \TriageSensB{}) and false alerts
(specificity \TriageSpecB{}), and they expressed higher confidence and spent roughly twice as long
per alert (mean confidence 4.2 vs.\ 2.2 in arm~A), consistent with reading a
richer explanation. In arm~A, near-chance balanced accuracy (\TriageBAA{}) confirms that the bare
anomaly score carries little actionable signal at the presented granularity. Inter-rater agreement
was higher in arm~B than arm~A (Cohen's $\kappa$: \TriageKappas), suggesting the explanation
creates a shared basis for judgment
that the score alone does not. By explanation family, arm-B accuracy was
\TriageFamIdn{} for identity, \TriageFamRes{} for resource, \TriageFamTmp{} for timing and
\TriageFamVol{} for volume explanations; performance was consistent across families, with
identity-deviation alerts reaching perfect accuracy (Table~\ref{tab:family}).

\begin{table}[t]
\centering
\caption{Triage study (H3): judgements pooled over four reviewers, 40 alerts (20 true).}
\label{tab:triage}
\resizebox{\columnwidth}{!}{%
\begin{tabular}{lcccccc}
\toprule
Arm & $n$ & Bal.\ acc. & Sens. & Spec. & Conf. & Sec. \\
\midrule
Score only & 80 & 0.475 & 0.62 & 0.33 & 2.2 & 19 \\
Explanation only & 80 & 0.913 & 0.90 & 0.93 & 4.2 & 40 \\

\bottomrule
\end{tabular}}
\end{table}

\begin{table}[t]
\centering
\caption{Precision of alerts by the behaviour family their explanation names (10 alerts/day, CERT
test windows, \ExpNAlerts{} alerts).}
\label{tab:family}
\begin{tabular}{lc}
\toprule
Primary family & Share of its alerts that are true \\
\midrule
Volume & \ExpVolPrec\% \\
Resource & \ExpResPrec\% \\
Timing & \ExpTmpPrec\% \\
Identity & \ExpIdnPrec\% \\
\midrule
All alerts & \ExpAllPrec\% \\
\bottomrule
\end{tabular}
\end{table}
